% !TeX root = ../教材题库_全部条目_转写工作稿.tex
% Source: XJ-B1-C05-PR-P195-Y0861
% 第5章 / 5.4 函数 y=A sin(ωx+φ) 的图象与性质 / 教材 p.195 / PDF 第202页
% OCR未校对：公式与文字均可能有误；未完成数学TeX转写；禁止直接发给学生
% 原题及相邻原解截图（核对用）：../题目原图/XJ-B1-C05-PR-P195-Y0861-01.webp
% 本文件只在转写完成后作为题干使用；OCR 条目尚不是可靠题干。
\begin{bankitem}{p.195 练习 第2题}{第5章 / 5.4 函数 y=A sin(\ensuremath{\omega}x+\ensuremath{\varphi}) 的图象与性质 / 教材 p.195 / PDF 第202页}{XJ-B1-C05-PR-P195-Y0861}{OCR未校对：公式与文字均可能有误；未完成数学TeX转写；禁止直接发给学生}
\begin{ocrdraft}
[教材p.195 / PDF第202页]
\par
2．填空：
\par
=sin，十《）的图象，只需把函数一sinx的图象上所
\par
（D为了得到函数3，
\par
有的点向
\par
平移
\par
个单位长度；
\par
=sin2，一一）的图象，只需把函数3，一sin2x的图象上
\par
（2）为了得到函数
\par
所有的点向
\par
平移
\par
个单位长度；
\par
（3）将函数3一、i！，，的图象上所有的点向右平移个单位长度，再把所得各
\par
点的横坐标伸长到原来的2倍（纵坐标不变），所得图象的函数解析式是
\end{ocrdraft}
\end{bankitem}
